\section{Saturation, variable quotients, and dual Koszul cocycles}
\label{app:saturation}

The relation-group theorem in Section~\ref{sec:depth} identifies
monomials after all prime variables are inverted. This appendix
examines what remains before localization: common multipliers,
quotients by one variable, and explicit dual Koszul classes.
Write $G=R/I$ for the binary ring and $S=R/J$ for its toric quotient.
At a finite variable cutoff put $\A_m=\{p\le m:p\text{ odd prime}\}$
and $w_m=\prod_{p\in\A_m}x_p$.

\subsection{Localization and saturation}

We ask whether the binary relations generate every
relation between monomials of the same bidegree. The toric map
\[
 G=R/I\twoheadrightarrow S=R/J
       =\kk[u^pv:p\text{ odd prime}],\qquad x_p\longmapsto u^pv,
\]
identifies all such monomials. By Proposition~\ref{setup:components},
$I=J$ precisely when every nonempty fiber is connected by binary
Goldbach moves, or equivalently when $\dim_\kk G_{n,d}\le1$ for every
$(n,d)$. The finite toric quotient $S^{<m}$ has dimension two for
$m>5$, so its depth is at most two. The full $S$ also has dimension
two; the associated projective curve has dimension one, as explained
in Section~\ref{sec:setup}. The surjection $G^{<m}\to S^{<m}$ alone,
however, does not give a depth bound for $G^{<m}$.

For a general finite support $\A$, put $w=\prod_{a\in\A}x_a$.
The toric map $G\twoheadrightarrow S$ from Section~\ref{sec:setup}
records only Goldbach degree and monomial degree. The group $U_{\A}$ retains the
binary relations and maps onto the subgroup of $\mathbb Z^2$ generated
by the bidegrees $(a,1)$, by $[a]\mapsto(a,1)$. In polynomial
coordinates the kernel of $G\to\kk[U_{\A}]$ is
$(I:w^\infty)/I$. In particular, two monomials $x^\alpha,x^\beta$
have the same image in this group algebra precisely when
$w^j x^\alpha=w^j x^\beta$ in $G$ for some $j\ge0$.
Here $I:w^\infty=\bigcup_{j\ge0}(I:w^j)$. In a finite primary
decomposition this saturation removes exactly the components whose
radicals contain at least one variable. Such a component may be
minimal or embedded. Since $S$ is a domain and $\phi(w)\ne0$,
\[
 I\subseteq (I:w)\subseteq (I:w^\infty)\subseteq J.
\]

Appendix~\ref{depth:prime-variable} computes these ideals at cutoff $42$
and explains how the colon enters the resolution of $G/x_p$.

For the full odd-prime support, Theorem~\ref{depth:prime-pair-group}
identifies $U_\infty$ with the subgroup of $\mathbb Z^2$
generated by the bidegrees $(p,1)$. Thus any two monomials
$x_{\boldsymbol p},x_{\boldsymbol q}$ of the same bidegree have the same
image after all variables are inverted. Clearing denominators gives a
finite monomial $x_{\boldsymbol r}$ such that
\[
 x_{\boldsymbol r}x_{\boldsymbol p}
     =x_{\boldsymbol r}x_{\boldsymbol q}\quad\text{in }G .
\]
Equivalently, $J=\bigcup_{\boldsymbol r}(I:x_{\boldsymbol r})$.
If the two monomials represent distinct classes in $G$, their
difference is a zero divisor, since $x_{\boldsymbol r}$ has nonzero
image in $S$.

In fact, in the full polynomial ring $R$ we have
\[
 J=\bigcup_m(I:w_m),\qquad w_m=\prod_{p\in\A_m}x_p .
\]
To see why a squarefree multiplier suffices, take $f\in J$ and a
monomial $x_{\boldsymbol r}$ with $x_{\boldsymbol r}f\in I$. For each
occurrence of a prime $p$ in $\boldsymbol r$ we can find distinct,
arbitrarily large primes $q,q'$ such that $q+q'-p$ is also prime.
Indeed, the three forms $a,b,a+b-p$ have pairwise nonproportional
linear parts. Modulo $2$ take $a=b=1$. At each odd prime $s$ take
$a=1$ and choose $b$ outside the two residues $0,p-1$. Thus the
linear forms theorem \cite[Theorem~1.6]{TT}, applied on
$[Y,2Y]\times[3Y,4Y]$ for large $Y$, gives the desired primes, with proper prime
powers discarded as in Lemma~\ref{depth:prime-pair-conversion}.
Choose each pair above all the preceding pairs. Since
$x_qx_{q'}=x_px_{q+q'-p}$ in $G$, the product of these distinct
variables is a multiple of $x_{\boldsymbol r}$ in $G$ and annihilates
the class of $f$. It divides $w_m$ for sufficiently large $m$. The
reverse inclusion follows from $I\subseteq J$ and $w_m\notin J$,
since $J$ is prime.

Here $m$ may depend on $f$. For the ideal of a fixed finite cutoff,
Theorem~\ref{depth:integral-stabilization} instead identifies the
relation group with the bidegree group once $m$ is sufficiently large.
Consequently $J^{<m}=I^{<m}:w_m^\infty$ for every such $m$.

\subsection{Quotients by a prime variable}\label{depth:prime-variable}
We use the variable cutoff on the odd primes. Throughout this
subsection $R=R^{<m}$, $I=I^{<m}$, $G=R/I$, and $S=R/J$ is the toric
quotient at the same cutoff. For $p\in\A_m$ write $G/x_p$ for
$G/x_pG$. All Betti numbers below are taken over $R$ unless another
polynomial ring is specified.

\paragraph{The cutoff $m=42$.}
Write $P_s=(x_p\mid p\le s,\ p\text{ odd prime})$ in $R^{<42}$.
Over $\mathbb Q$ and $\mathbb F_2$, a primary decomposition gives
\[
 I^{<42}=J^{<42}\cap P_{29}\cap Q,\qquad \sqrt Q=P_{31},
\]
where $J^{<42}$ is the kernel of $x_p\mapsto u^pv$ on $R^{<42}$.
One choice of the embedded component is
$Q=(I^{<42}+(x_{31}^3)):x_{37}^\infty$.
The two minimal primes are $J^{<42}$ and $P_{29}$.
Thus the variable zero divisors are exactly the $x_p$ with $p\le31$.
Both $x_{37}$ and $x_{41}$ are regular. The supplement verifies the
decomposition and the variable colon ideals in \textsc{Singular}.
Here one colon suffices. The calculation over $\mathbb Q$ gives
$I^{<42}:x_3=J^{<42}$. This equality holds over every field: for each
binomial of $J^{<42}$, membership of its $x_3$-multiple in $I^{<42}$
means that its two monomials after multiplication are connected by
binary moves (Proposition~\ref{setup:components}). Such binomials
span $J^{<42}$ over each field. Hence also $I^{<42}:w_{42}=J^{<42}$.
The zero-divisor variables need not be an initial segment at other
cutoffs. Over the same fields, $x_{19}$ is regular in $G^{<30}$
whereas $x_{23}$ is a zero divisor.

One can see the loss of information under localization directly over
every field. Take $f=x_{13}x_{37}^2-x_5x_{41}^2$. Its monomials both
map to $u^{87}v^3$ in the toric ring. At cutoff $42$ the component of
$x_{13}x_{37}^2$ consists of this monomial and $x_{19}x_{31}x_{37}$,
so $f\ne0$ in $G^{<42}$. Over the two fields above the colon ideal
$(I^{<42}:f)$ is exactly $P_{31}$. One can check $x_3 f=0$ directly
by the three moves
\[
 x_3x_{13}x_{37}^2=x_5x_{11}x_{37}^2
 =x_5x_7x_{37}x_{41}=x_3x_5x_{41}^2 .
\]
At cutoff $43$ the new prime $43$ connects the original monomials by the path
\[
 x_{13}x_{37}^2=x_7x_{37}x_{43}=x_3x_{41}x_{43}=x_5x_{41}^2 .
\]
Thus $f$ dies on passing to $G^{<43}$, the first variable cutoff
at which $43$ is available.

\begin{proposition}[Betti numbers after setting a variable to zero]
\label{depth:variable-betti}
Fix $p\in\A_m$ and let $I'=I:x_p$. For $i\ge1$,
\[
 \beta_{i,n,i+1}(G/x_p)
 =\beta_{i,n,i+1}(G)+\beta_{i-1,n-p,i}(S).
\]
If $\operatorname{reg}_R G\le2$, then in every bidegree and for $i\ge0$,
\begin{equation}\label{depth:quotient-betti}
 \beta_{i,n,d}(G/x_p)
 =\beta_{i,n,d}(G)+\beta_{i-1,n-p,d-1}(R/I'),
\end{equation}
where a Betti number in negative homological degree is zero.
In particular the last term is a Betti number of $S$ when $I'=J$.
\end{proposition}
\begin{proof}
Since $I$ contains all quadratic equal-sum binomials on $\A_m$,
the map $G\to S$ is an isomorphism in monomial degrees $0,1,2$.
Also $I\subseteq I'\subseteq J$, because the image of $x_p$ in $S$ is
nonzero. Thus $G\to R/I'\to S$ identifies the linear strands in Tor:
in monomial degree $i+1$, the three Koszul terms computing $\Tor_i$ have
coefficient monomial degrees $0,1,2$. The same comparison applies to
$G/x_p\to S/x_pS$, which again agrees through monomial degree two.
Multiplication by $x_p$ is injective on $S$. The mapping cone of
multiplication by $x_p$ on its minimal resolution is therefore a
minimal resolution of $S/x_pS$, giving the first formula.

For the second formula use the exact sequence
\[
 0\longrightarrow (R/I')(-p,-1)
 \xrightarrow{\ \cdot x_p\ } G\longrightarrow G/x_p\longrightarrow0.
\]
We show that the first map induces zero on Tor. In homological degree
zero this follows from $x_p\in\m$. For $i\ge1$, the first possible
monomial degree in $\Tor_i^R(R/I',\kk)$ is $i+1$, since $I'$ has no linear forms.
In that monomial degree the classes lift from $\Tor_i^R(G,\kk)$ by the linear
strand comparison. The composite $G(-p,-1)\to(R/I')(-p,-1)\to G$
is multiplication by $x_p$, which acts as zero on Tor with $\kk$.
Classes of greater monomial degree have shifted target monomial degree at least $i+3$,
where $\Tor_i^R(G,\kk)$ vanishes by the regularity bound.
The long exact sequence now gives \eqref{depth:quotient-betti}.
\end{proof}

At cutoff $42$, over both $\mathbb Q$ and $\mathbb F_2$, the supplement
verifies $\operatorname{reg}_R G\le2$ and gives
\[
\begin{array}{c|ccc}
 &G&S&G/x_3\\ \hline
 \operatorname{depth}_R &2&1&0\\
 \pd_R &10&11&12
\end{array}
\]
Thus \eqref{depth:quotient-betti} applies with $p=3$ and $R/I'=S$. For
example, $x_5x_7$ is a nonzero class in $G/x_3$ killed by every
variable. It is one of six socle classes, represented by
\[
 x_5x_7,\quad x_5x_{13},\quad x_5x_{19},\quad
 x_5x_{23},\quad x_7x_{23},\quad x_5x_{31}.
\]
Since $\pd_R G=10$, the connecting map identifies these six top
Koszul classes in $\Tor_{12}^R(G/x_3,\kk)_{n,d}$ with
$\Tor_{11}^R(S,\kk)_{n-3,d-1}$, bidegree by bidegree.

\paragraph{Powers of \texorpdfstring{$x_3$}{x3}.}
For a fixed finite cutoff the ascending chain $(I:x_3^j)$ stabilizes
by Noetherianity. There is a simple criterion for its limit to be $J$.
For $m>5$ this happens precisely when every prime $7\le p\le m$ admits
a representation $p+3=q+q'$ with $q,q'>3$ prime. Indeed, $q,q'<p$,
so after inverting $x_3$ the relations $x_3x_p=x_qx_{q'}$ express
each $x_p$ recursively in $x_3^{\pm1},x_5$. The toric images
$u^3v,u^5v$ are algebraically independent, so the resulting localized
map to the localized toric ring is injective. Contracting to $R$ gives
$I:x_3^\infty=J$.

Conversely, if $p$ has no such representation, put $j=(p-3)/2$.
The binomial $x_3^{j-1}x_p-x_5^j$ belongs to $J$. After multiplication
by any power of $x_3$, the first monomial remains isolated under binary
moves: each available pair has sum $6$ or $3+p$, both uniquely represented
on this support. Consequently no power of $x_3$ times this binomial belongs to $I$.
A failure of the pair condition at $p$ therefore persists at every
cutoff containing $p$. It follows that $I^{<m}:x_3^\infty=J^{<m}$ for
arbitrarily large $m$ precisely when every odd prime $p\ge7$ meets the
pair condition. We have no general bound on the power needed when
equality holds, and the results here leave open whether
$I^{<m}:x_3=J^{<m}$ for infinitely many $m$. The equality
$J=\bigcup_m(I:w_m)$ for the full ring allows the multiplier to
depend on the binomial.

\begin{example}[The first failure of a single colon]\label{depth:cutoff132}
At every nonempty variable cutoff before $m=131$ we have $I:x_3=J$.
At $m=131$, set
\[
 f=x_{127}^6-x_{107}x_{131}^5,\qquad I'=I:x_3.
\]
The two monomials have bidegree $(762,6)$, and
\[
 f\notin I',\qquad J=I'+(f)=I:x_3^2.
\]

To see why one copy of $x_3$ fails, consider the two partitions of
$765$ obtained after multiplying by $x_3$. Starting on the left, two legal moves give
\[
 x_3x_{127}^6 \longrightarrow x_{23}x_{107}x_{127}^5
 \longrightarrow x_{19}x_{107}x_{127}^4x_{131}.
\]
Sort the primes in increasing order and then reduce each entry modulo $3$.
The two moves give the residue words
\[
 (0,1,1,1,1,1,1)\ \longrightarrow\
 (2,2,1,1,1,1,1)\ \longrightarrow\
 (1,2,1,1,1,1,2).
\]
The first move, $3+127=23+107$, replaces $(0,1)$ by $(2,2)$ and
uses the only $3$. The second, $23+127=19+131$, preserves the
residue counts even though it changes the primes. To turn two entries
congruent to $1$ back into a $3$ and an entry congruent to $2$ at
this vertex, the only candidate sums are $19+127=3+143$ and
$127+127=3+251$. The first requires the composite $143$ and the
second requires $251>131$. The two entries congruent to $2$ would need
$107+131=3+235$, with $235=5\cdot47$. Thus no next move from this
vertex can reintroduce $3$. Here the residue counts constrain which moves are possible. To follow a
loop as a Koszul cycle, Example~\ref{ex:weight33} also keeps track of
which prime is fixed on each edge and the order of the two fixed primes
at each corner.

Every partition in this fiber has seven entries and at most one $3$:
two copies of $3$ would give a sum at most $6+5\cdot131=661$.
Write $h$ for the count of entries congruent to $1$ and $z$ for the
number of threes. Reduction of the sum modulo $3$ gives
$2h-(7-z)\equiv0\pmod3$. Thus
\[
 \begin{array}{c|c}
 z & h\text{ allowed by the congruence}\\ \hline
 0 & 2,\ 5\\
 1 & 0,\ 3,\ 6
 \end{array}
\]
A move replaces only two entries, so it changes $h$ by at most two.
The two endpoints have $h=6$ and $h=0$, so any path between them must pass through $h=3$. Such a partition
would contain one $3$ and \emph{six other entries}, three from each
nonzero residue class. These six entries must sum to $765-3=762$.
Their largest possible sum is $3\cdot127+3\cdot131=774$. The
next-largest primes in the two classes are $109$ and $113$, respectively.
Any reduction therefore costs at least $18$, whereas the required
reduction is $12$. This excludes the intermediate partition and
proves that $x_3f\notin I$.

A second $3$ allows us to continue after the same two moves from the
left. Use the remaining $3$ in $3+19=11+11$, and then use
$11+127=7+131$ and $7+127=3+131$ twice in succession. This gives
$x_3^2x_{107}x_{131}^5$, hence $x_3^2f\in I$. The obstruction is
specific to the cutoff. When $137$ enters, a path with just one $3$
can pass through $(3,113,127,127,127,131,137)$. Its residue word
$(0,2,1,1,1,2,2)$ has $h=3$, with all three entries congruent to
$1$ coming from $127$.

For completeness, the supplement checks each odd-prime prefix through
$131$. It verifies the pair condition above and uses \textsc{Singular}
over $\mathbb Q$ to compare consecutive colons. They agree at the first
power for each earlier prefix. At the prefix ending at $131$, the second and third colons agree,
and the second is $I'+(f)$. These equalities hold over every field.
Indeed, membership of a binomial in $I:x_3^j$ is connectivity of its
endpoints after multiplication by $x_3^j$. These colons are spanned by
such binomials, by the component basis. Adjoining $f$ identifies the
components joined by its monomial multiples, again independently of
the field.

The calculation also gives $I':f=(x_p\mid p\le113)$ at this cutoff.
Each of those variables kills the class of $f$ over every field by
the same connectivity test. Every $x_{127}^a x_{131}^b f$ survives,
and distinct pairs $(a,b)$ give different bidegrees. It follows that
\[
 J/I'\cong\kk[x_{127},x_{131}](-762,-6).
\]
This describes the difference between the actual term $R/I'$ of the
colon sequence and the toric quotient $S$. The Koszul complex on the
$29$ variables through $113$ resolves $J/I'$, with the indicated shift.
Consequently $0\to J/I'\to R/I'\to S\to0$ induces isomorphisms
$\Tor_i^R(R/I',\kk)_{n,d}\cong\Tor_i^R(S,\kk)_{n,d}$ outside the two
monomial degrees $i+5$ and $i+6$.
\end{example}

\subsection{Dual Koszul cocycles and the prime-pair condition}
\label{depth:dual-cocycles}

The depth bound can also be approached by constructing classes in the
dual Koszul complex. At a variable cutoff write $R_m=R^{<m}$ and
$G_m=G^{<m}$, and let $\mathbb K_{\kk}^{<m}$ be the Koszul resolution of
$\kk$ over $R_m$. Put
\[
 C_m=\operatorname{Hom}_{R_m}(\mathbb K_{\kk}^{<m},G_m).
\]
Write $C_{m,n,d}$ for its complex in internal bidegree $(n,d)$ and
$C^i_{m,n,d}$ for the term in cohomological degree $i$.
In this subsection we use $\varepsilon_q$ for the dual exterior
generators, of cohomological degree one and internal bidegree
$(-q,-1)$. With the transpose convention for the codifferential,
\[
 \mathrm{d}^\star\omega=\theta_m\wedge\omega,
 \qquad \theta_m=\sum_{q\in\A_m}x_q\varepsilon_q.
\]
Thus $\mathrm{d}^\star$ preserves internal bidegree, and
$\mathrm H^i(C_m)=\operatorname{Ext}_{R_m}^i(\kk,G_m)$.
The least $i$ with nonzero cohomology is $\delta^{<m}$.
If $s_m=\sum_{q\in\A_m}q$, finite Koszul self-duality gives
\begin{equation}\label{depth:dual-betti}
 \dim_\kk\mathrm H^i(C_m)_{n,d}
 =\beta^{<m}_{|\A_m|-i,\,s_m+n,\,|\A_m|+d}.
\end{equation}

Suppose $m\ge11$ and every prime $5\le q\le m$ has a representation
$q+1=p_1+p_2$ by odd primes. Both summands are at most $q-2$, so they
belong to $\A_m$. Choose $p_1\le p_2$ with $p_1p_2$ maximal, using
the product order of Proposition~\ref{moves:order}, and set
\[
 y_q=x_{p_1}x_{p_2}\quad\text{in }G_m.
\]
The maximum-product choice fixes a representative; the class $y_q$
is independent of this choice by the defining quadratic relations.
In particular,
\[
 y_5=x_3^2,\qquad y_7=x_3x_5,\qquad y_{11}=x_5x_7.
\]
Consider the cochain
\begin{equation}\label{depth:dual-family}
 \omega_m=\varepsilon_3\wedge
       \sum_{\substack{5\le q\le m\\q\text{ prime}}}y_q\varepsilon_q
       \in C^2_{m,-2,0}.
\end{equation}

\begin{proposition}\label{depth:dual-family-criterion}
The cochain $\omega_m$ is closed if and only if
\begin{equation}\label{depth:dual-cubic-identities}
 x_qy_{q'}=x_{q'}y_q\quad\text{in }G_m
 \qquad(5\le q<q'\le m).
\end{equation}
Equivalently, if $q+1=p_1+p_2$ and $q'+1=p'_1+p'_2$ are the
chosen representations, then
\[
 x_q(x_{p'_1}x_{p'_2})=x_{q'}(x_{p_1}x_{p_2})\quad\text{in }G_m.
\]
Whenever it is closed, its class is nonzero over every field. In
particular,
\[
 \beta^{<m}_{|\A_m|-2,\,s_m-2,\,|\A_m|}\ge1,
 \qquad \delta^{<m}\le2.
\]
\end{proposition}
\begin{proof}
Exterior multiplication gives
\[
 \mathrm{d}^\star\omega_m=
 \sum_{5\le q<q'\le m}(x_{q'}y_q-x_qy_{q'})
        \varepsilon_3\varepsilon_q\varepsilon_{q'},
\]
which proves the criterion. A possible primitive in internal
bidegree $(-2,0)$ has the form $\omega'=\sum_q f_q\varepsilon_q$, with
$f_q\in(G_m)_{q-2,1}$. Since $1$ and $9$ are not prime labels,
$f_3=f_{11}=0$. The coefficient of
$\varepsilon_3\varepsilon_{11}$ in every such $\mathrm{d}^\star\omega'$
is therefore $x_3f_{11}-x_{11}f_3=0$. In $\omega_m$ that coefficient is the nonzero
monomial class $x_5x_7$. The component basis of
Proposition~\ref{setup:components} proves its nonvanishing over every
field. Equation~\eqref{depth:dual-betti} and the depth characterization
give the final assertions.
\end{proof}

\begin{example}[Finite cocycles]\label{depth:dual-finite}
At cutoff $12$ the cocycle is
\[
 \omega_{12}=x_3^2\varepsilon_3\varepsilon_5
       +x_3x_5\varepsilon_3\varepsilon_7
       +x_5x_7\varepsilon_3\varepsilon_{11}.
\]
Its closure follows from $x_3x_7=x_5^2$ and $x_3x_{11}=x_7^2$.
For example,
$x_5y_{11}=x_5^2x_7=x_3x_7^2=x_3^2x_{11}=x_{11}y_5$.
The following choices extend the displayed expression by one term
whenever a new prime is adjoined:
\[
\begin{array}{c|ccccc}
 q&13&17&19&23&29\\ \hline
 y_q&x_7^2&x_7x_{11}&x_7x_{13}&x_{11}x_{13}&x_{13}x_{17}
\end{array}
\]
Thus $\omega_{14}=\omega_{12}+x_7^2\varepsilon_3\varepsilon_{13}$ and
$\omega_{18}=\omega_{14}+x_7x_{11}\varepsilon_3\varepsilon_{17}$, and
similarly at cutoffs $20$, $24$, and $30$.

The supplement verifies~\eqref{depth:dual-cubic-identities} at every
odd-prime prefix through $101$. Hence $\omega_m$ is a nonboundary cocycle
over every field for every integer $11\le m\le102$.
The verification enumerates cubic monomials and their binary-move
components at each support; it does not replace $G_m$ by its toric
quotient. The coefficient test in the preceding proof certifies
nontriviality independently of the field.

At the smaller cutoff $8$ there is instead a linear-coefficient class
\[
 \omega_8=x_3\varepsilon_3\varepsilon_5+x_5\varepsilon_3\varepsilon_7
       \in C^2_{8,-5,-1}.
\]
Its codifferential is
$\mathrm{d}^\star\omega_8=(x_3x_7-x_5^2)\varepsilon_3\varepsilon_5\varepsilon_7=0$.
The only possible primitive in this bidegree is a scalar multiple
of $\varepsilon_5$; the codifferential of every nonzero such multiple
has a nonzero $\varepsilon_5\varepsilon_7$ coefficient, so the displayed class is
nonzero.
\end{example}

\begin{conjecture}[Infinite dual Koszul cocycle]\label{depth:dual-conjecture}
For every odd prime $q\ge5$, there are odd primes $p_1,p_2$ with
$q+1=p_1+p_2$. Let $y_q=x_{p_1}x_{p_2}\in G_{q+1,2}$ be the
resulting quadratic class in the full binary Goldbach ring. Then
\[
 x_qy_{q'}=x_{q'}y_q\quad\text{in }G\qquad(5\le q<q').
\]
\end{conjecture}

For infinitely many variables,
$C^i=\operatorname{Hom}_R(\mathbb K_{\kk,i},G)$ is a product over
the $i$-element exterior supports. Each coordinate of $\mathrm{d}^\star$ uses
only finitely many input coordinates. Under the conjecture,
\[
 \omega=\varepsilon_3\wedge\sum_{q\ge5}y_q\varepsilon_q
\]
therefore defines a cocycle in $C^2$. It also belongs to the
completion of the direct-sum cochain module in the Goldbach weight
of its coefficients: the coefficient $y_q$ has weight $q+1$.
This is a coefficient filtration; every whole term has dual internal
bidegree $(-2,0)$. The same $\varepsilon_3\varepsilon_{11}$
coefficient test proves that $\omega$ is not a boundary, even in this
completion. In particular, the conjecture gives
$\operatorname{Ext}_R^2(\kk,G)_{-2,0}\ne0$.

The conjecture concerns identities in the full ring. Requiring every
$\omega_m$ to be closed already in $G^{<m}$ is stronger; the finite
examples have this additional property. A cubic identity of weight
$q+q'+1$ may in general use primes as large as $q+q'-5$ along a move
path, beyond the original labels $q,q'$.

\begin{proposition}\label{depth:dual-prime-pair}
Conjecture~\ref{depth:dual-conjecture} implies that every prime
$q\ge7$ has a representation
\[
 q+3=p_1+p_2,\qquad p_1,p_2>3\text{ prime}.
\]
Consequently it implies $I^{<m}:x_3^\infty=J^{<m}$ at every
variable cutoff $m>5$.
\end{proposition}
\begin{proof}
The identity with the label $5$ is $x_5y_q=x_3^2x_q$.
If $q+3$ has no representation by primes greater than $3$, the
monomial $x_3^2x_q$ is isolated under binary moves: its only pair
sums are $6$ and $q+3$, each uniquely represented. For $q\ge7$ the
monomial $x_5y_q$ is different, since it contains $x_5$.
The component basis therefore contradicts the required equality.
Moreover, every prime occurring along a cubic path of weight $q+6$
is at most $q$, so larger cutoffs cannot repair a failure of this
particular identity. The saturation conclusion follows from the
prime-pair criterion in Appendix~\ref{depth:prime-variable}.
\end{proof}

The relation-group theorem does not supply the cubic identities in
Conjecture~\ref{depth:dual-conjecture}. When $y_q,y_{q'}$ exist,
$x_qy_{q'}-x_{q'}y_q$ belongs to $J$, so that theorem supplies a common
monomial multiplier annihilating the difference in $G$.
It does not permit cancellation of that multiplier. Even saturation
by powers of $x_3$ only gives annihilation by some power, rather
than equality in $G$. Thus the conjecture asks for specific cubic
binary-move paths beyond the stable connectivity already proved.

